%=============================================================================
% WAVE 4, ITERATION 13: EXPERIMENTAL ROADMAP UPDATE
% Post-Dark-SRF Recalibration + Six New Experimental Proposals
%
% Agent 17 (Experimental Proposals) | DFD Research Programme | Model: Opus 4.6
% Building on: W4i10_Experimental_Proposals.tex, W4i12_DarkSRF_Dictionary.tex,
%              MASTER_FINDINGS_CORPUS.md, section02_formalism.tex
%
% Status flags: [VERIFIED], [NEW], [ACTIONABLE], [CORRECTED]
% Date: 20 March 2026
%=============================================================================

\documentclass[11pt]{article}
\usepackage[utf8]{inputenc}
\usepackage{amsmath,amssymb,amsfonts,mathtools}
\usepackage{geometry}
\geometry{margin=1in}
\usepackage{booktabs}
\usepackage{siunitx}
\usepackage{hyperref}
\usepackage{xcolor}
\usepackage{enumitem}
\usepackage{float}
\usepackage{array}
\usepackage{longtable}
\usepackage{tcolorbox}
\tcbuselibrary{skins,breakable}

\newcommand{\ee}[1]{\times 10^{#1}}
\newcommand{\lbone}{|\lambda - 1|}
\newcommand{\dpsi}{\delta\psi}
\newcommand{\Qpsi}{Q_\psi}
\newcommand{\SNR}{\mathrm{SNR}}
\newcommand{\LCDM}{\Lambda\text{CDM}}

\definecolor{proposalcolor}{rgb}{0.0,0.35,0.55}
\definecolor{actioncolor}{rgb}{0.0,0.5,0.0}
\definecolor{newcolor}{rgb}{0.6,0.0,0.6}
\definecolor{warncolor}{rgb}{0.8,0.2,0.0}

\newcommand{\confirmed}[1]{\textcolor{blue}{\textbf{CONFIRMED:} #1}}
\newcommand{\corrected}[1]{\textcolor{red}{\textbf{CORRECTED:} #1}}
\newcommand{\newfinding}[1]{\textcolor{teal}{\textbf{NEW FINDING:} #1}}
\newcommand{\caution}[1]{\textcolor{orange}{\textbf{CAUTION:} #1}}

\title{Wave 4, Iteration 13: Experimental Roadmap Update\\
\large Post-Dark-SRF Recalibration $\cdot$ Six New Proposals $\cdot$ Optimal Sequencing\\[6pt]
\normalsize Density Field Dynamics Research Programme}

\author{DFD Research Programme --- Experimental Proposals Agent (W4i13, Opus 4.6)\\
\textit{Building on: W4i10 Experimental Proposals, W4i12 Dark SRF Dictionary,\\
MASTER\_FINDINGS\_CORPUS, section02\_formalism}}

\date{20 March 2026}

\begin{document}
\maketitle

%=============================================================================
\section*{Executive Summary: Top 5 Findings}
%=============================================================================

\begin{tcolorbox}[colback=blue!5!white, colframe=blue!60!black,
  title=\textbf{W4i13 TOP 5 FINDINGS --- Ranked by Programme Impact}]
\begin{enumerate}

\item \textbf{FINDING 1 [NEW, ACTIONABLE]: DESI DR2 $D_H/r_d$ sign-change test
is the single highest-impact experiment and costs \$0.} DFD predicts a
non-monotonic residual in $D_H/r_d$ relative to $\LCDM$ with a sign
change at $z = 1.78 \pm 0.12$. DESI DR2 (March 2025, 14M+ tracers) covers
$0.1 < z < 4.2$ with measured $D_H/r_d$ at 7 effective redshift bins.
CPL models predict monotonic departure; DFD predicts a turnover.
\textbf{This is testable NOW with published data} --- no new observations
needed. A single competent cosmologist with the DESI DR2 public data
release can produce the residual plot in 1--2 weeks.
\textbf{Derivation chain}: $\psi$-screen optical bias (section02, Eq.~2.1)
$\to$ $n = e^\psi$ modifies $D_A$ and $D_H$ $\to$ $\mu$-crossover at
$a_0/a(z)$ produces non-monotonic effective equation of state $\to$
$D_H/r_d$ residual sign change at $z_{\rm cross} = 1.78 \pm 0.12$
(W4i11, Prediction D4).

\item \textbf{FINDING 2 [NEW]: Milligram-scale torsion pendulums are now
viable DFD test platforms.} Two 2025--2026 breakthroughs:
(a)~1~mg SiN$_x$ torsion pendulum cooled to 240~$\mu$K at 18~Hz
(Communications Physics 2026), surpassing prior cooling records by
$20\times$; (b)~nanofabricated Si$_3$N$_4$ torsion fiber supporting
87~g test masses with atto-Newton force sensitivity
(arXiv:2601.11366, January 2026). These platforms can test DFD's
predicted scalar-gradient force $F_\psi = (c^2/2)\,m\,\nabla\psi$
near a controlled EM source mass (high-$Q$ SRF cavity as psi-source).
At $Q_{\rm EM} = 5\ee{10}$ and stored energy $U = 10$~J, the predicted
$\dpsi_{\rm EM}$ gradient produces a force $\sim 10^{-19}$~N at 1~cm,
which is within $10\times$ of demonstrated atto-Newton sensitivity.
\textbf{This opens a new detection channel not in W4i10.}

\item \textbf{FINDING 3 [CORRECTED]: Post-Dark-SRF experimental roadmap
resequenced. SQMS Phase~I is confirmed as gateway, but the D$_H$/r$_d$
test now ranks above it on cost-effectiveness.}
With Dark SRF INVALIDATED (W4i12), the SQMS bound $\lbone < 1.56\ee{-9}$
is authoritative, and the Phase~I discovery window is \emph{wider} than
previously estimated. The recalibrated optimal sequence is:
(0)~$D_H/r_d$ sign-change test [\$0, 2 weeks];
(0b)~GPS 59~ps archival [\$50--100K, 6 months];
(0c)~LIGO O4 scalar reanalysis [\$100--500K, 12 months];
(I)~SQMS Phase~I [\$3.2M, 24 months];
(II)~Si$_3$N$_4$ MEMS prototype [\$200--500K, 12--18 months];
then torsion-pendulum and optical-clock proposals.
\textbf{Derivation}: Dark SRF produced a claimed bound of
$4.2\ee{-10}$--$7.1\ee{-10}$ (W4i10/W4i11), which would have narrowed
the Phase~I discovery window. W4i12 shows the EM-to-$\psi$ conversion
uses gravitational coupling $G/c^4 = 9.3\ee{-45}$, making active LSW
hopelessly weak (Passive Superiority Theorem, W3i5). With this
invalidation, the entire $1.56\ee{-9}$ to $1.9\ee{-10}$ range is
now open for Phase~I discovery.

\item \textbf{FINDING 4 [NEW]: Optical clock altitude comparison can
probe DFD's scalar-field-dependent redshift at the $10^{-19}$ level.}
Transportable optical lattice clocks have achieved $10^{-18}$ fractional
frequency comparison (Tokyo Skytree, 450~m height, 2023), and CU~Boulder
is deploying clocks to Mt.~Blue Sky (4,350~m elevation) in 2025--2026
via free-space laser link. DFD predicts the gravitational redshift is
$\Delta f/f = \Delta\psi/2 + (\lambda-1)\,\Delta\psi_{\rm EM}/(2c^2)$,
with the second term a \emph{nonreciprocal} correction absent in GR.
For the Sun-sourced $\psi_{\rm EM}$ gradient between sea level and
4,350~m, $\Delta\psi_{\rm EM} \approx 1.8\ee{-17}$.
At $\lbone \sim 10^{-9}$ (SQMS bound), the anomalous redshift correction
is $\sim 10^{-26}$ --- far below clock precision. \textbf{However}, the
\emph{diurnal modulation} of the EM-sourced term (tracking the Sun's
position) provides a distinctive temporal signature that can be
integrated over months to improve sensitivity by $\sqrt{N_{\rm days}}$.
With 1 year of continuous comparison at $10^{-19}$/comparison,
the integrated sensitivity reaches $\sim 10^{-20}$, still 6 orders
below the predicted signal. \textbf{Verdict: not yet competitive with
SQMS, but establishes methodology for future clock networks
(CLONETS-DS, 2028).}

\item \textbf{FINDING 5 [NEW]: Neutron interferometry with a scalar-field
source provides an independent test geometry distinct from all other
DFD experiments.}
The screened scalar field literature (arXiv:2405.14638, 2024) shows
that neutron interferometers are uniquely sensitive to environmental
scalar field gradients because neutrons couple to $\psi$ through their
rest-mass energy. A DFD-specific neutron interferometry experiment
would place a high-$Q$ SRF cavity (psi-source) near one arm of a
perfect-crystal Si interferometer at a facility like NIST or ILL.
The predicted phase shift is
$\Delta\phi = m_n c^2 \dpsi_{\rm EM} L / (\hbar v_n)$, where $L$
is the path length and $v_n$ is the neutron velocity. For thermal
neutrons ($v_n = 2200$~m/s, $\lambda_n = 1.8$~\AA), $L = 5$~cm, and
$\dpsi_{\rm EM} \sim 10^{-36}$ (from an SRF cavity with $Q = 5\ee{10}$,
$U = 10$~J), the predicted phase shift is
$\Delta\phi \sim 10^{-22}$~rad --- far below measurable ($\sim 10^{-4}$~rad).
\textbf{Verdict: KILLED. The gravitational coupling $G/c^4$ makes
EM-sourced $\dpsi$ too small for neutron interferometry.} This is
another manifestation of the Passive Superiority Theorem: active
EM-sourced psi-signals are generically too weak.

\end{enumerate}
\end{tcolorbox}

\tableofcontents
\newpage


%=============================================================================
%=============================================================================
\part{Post-Dark-SRF Roadmap Recalibration}
\label{part:roadmap}
%=============================================================================
%=============================================================================

%=============================================================================
\section{What Changed with Dark SRF Invalidation (W4i12)}
%=============================================================================

The W4i12 Dark SRF Dictionary established three critical facts that
force a recalibration of the experimental roadmap:

\begin{enumerate}
\item \textbf{The DFD coupling uses gravitational strength $G/c^4 =
9.3\ee{-45}$~N$^{-1}$}, making active EM-to-$\psi$ light-shining-through-wall
conversion hopelessly weak. The Dark SRF claimed bounds of $\lbone <
4.2\ee{-10}$ (W4i10) and $7.1\ee{-10}$ (W4i11) are both \textbf{incorrect}
--- they conflated DFD's spin-0 scalar coupling with the dark photon's
spin-1 kinetic mixing. The two have fundamentally different physics:
dark photons couple via $\chi \mathbf{E}\cdot\mathbf{E}'$ (vector-vector),
while DFD's $\dpsi$ is sourced by $u_{\rm EM} \propto (E^2 + B^2)/2$
(scalar total energy density) with strength $(\lambda-1) \cdot 4\pi G/c^4$.

\item \textbf{The SQMS bound $\lbone < 1.56\ee{-9}$ is restored as
authoritative.} This bound comes from the Q-factor of SRF cavities
(passive measurement), not from active LSW conversion. It is immune
to the Dark SRF invalidation because it measures a different observable:
the psi-field's back-reaction on EM cavity losses, not forward conversion.

\item \textbf{The SQMS Phase~I discovery window is WIDER} than W4i10/W4i11
estimated. With the Dark SRF ``bound'' removed, the entire range
$1.56\ee{-9}$ down to the Phase~I target of $1.9\ee{-10}$ (at $Q_0 = 2\ee{11}$)
is open territory. No other experiment constrains this range.
\end{enumerate}

\subsection{Impact on W4i10 proposals}

\begin{center}
\renewcommand{\arraystretch}{1.3}
\begin{tabular}{@{}clcc@{}}
\toprule
\textbf{\#} & \textbf{W4i10 Proposal} & \textbf{W4i12 Impact} & \textbf{W4i13 Status} \\
\midrule
1 & LIGO O3/O4 Archival & None (passive) & \textbf{RETAINED} \\
2 & GPS 59~ps Archival & None (passive) & \textbf{RETAINED} \\
3 & Si$_3$N$_4$ MEMS Prototype & None (passive) & \textbf{RETAINED} \\
4 & Galaxy Rotation Curve & None (astrophysical) & \textbf{RETAINED} \\
\bottomrule
\end{tabular}
\end{center}

All four W4i10 proposals survive the Dark SRF invalidation because they are
all passive observations or astrophysical data analyses. The Dark SRF
invalidation \emph{only} kills active EM-to-$\psi$ conversion experiments.
The Passive Superiority Theorem (W3i5) is validated yet again.


%=============================================================================
\section{Recalibrated Detection Ranking}
\label{sec:ranking}
%=============================================================================

\begin{tcolorbox}[colback=green!3!white, colframe=actioncolor,
  title=\textbf{UPDATED EXPERIMENTAL RANKING (Post-W4i12)}]

\begin{center}
\renewcommand{\arraystretch}{1.3}
\begin{tabular}{@{}clccl@{}}
\toprule
\textbf{Rank} & \textbf{Experiment} & \textbf{Cost} & \textbf{Timeline} &
  \textbf{DFD Reach / Observable} \\
\midrule
\rowcolor{green!10}
0a & $D_H/r_d$ sign change (DESI DR2) & \$0 & 2 weeks &
  Sign change at $z = 1.78 \pm 0.12$ \\
\rowcolor{green!10}
0b & GPS 59~ps archival & \$50--100K & 6 months &
  SNR $> 100$; ratio $R = 1.070$ \\
\rowcolor{green!10}
0c & LIGO O4 scalar reanalysis & \$100--500K & 12 months &
  $\lbone \sim 1.5\ee{-14}$ \\
\midrule
I & SQMS Phase~I (4-cavity) & \$3.2M & 24 months &
  $\lbone$: $7.8\ee{-10}$ to $1.9\ee{-10}$ \\
\midrule
II & Si$_3$N$_4$ MEMS prototype & \$200--500K & 12--18 months &
  $Q_\psi$ threshold $2.2\ee{8}$ \\
\midrule
\rowcolor{purple!5}
N1 & Torsion pendulum scalar force & \$300--600K & 18--24 months &
  Force $\sim 10^{-19}$~N (near-term) \\
\rowcolor{purple!5}
N2 & Optical clock altitude modulation & \$0 (piggyback) & 12 months &
  Methodology for CLONETS-DS \\
\bottomrule
\end{tabular}
\end{center}

Rows marked green are zero-cost or archival experiments.
Rows marked purple are \textbf{NEW proposals from W4i13}.
\end{tcolorbox}

\textbf{Key change from W4i10}: The $D_H/r_d$ sign-change test is
elevated to Rank~0a (above all other experiments) because:
\begin{itemize}[nosep]
\item It costs \$0 and takes 2 weeks.
\item It tests a \emph{zero-parameter} DFD prediction (the sign change
  location depends only on the $\mu$-crossover, not on $\lambda$).
\item A positive detection (non-monotonic $D_H/r_d$ residuals with
  turnover near $z \sim 1.8$) would be strong evidence for DFD over
  CPL dark energy models.
\item A negative result (monotonic residuals) would sharpen the DESI
  tension from 3.1--3.5$\sigma$ to $>4\sigma$, approaching falsification
  of DFD's cosmological sector.
\item DESI DR2 data is \emph{already published} (March 2025).
\end{itemize}


%=============================================================================
\section{Optimal Experimental Sequence}
\label{sec:sequence}
%=============================================================================

The post-Dark-SRF optimal sequence accounts for three constraints:
(i)~cost-effectiveness (cheapest decisive tests first),
(ii)~information flow (each result informs the next investment),
(iii)~personnel availability (parallel tracks where possible).

\subsection{Phase 0: Zero-Cost and Archival (NOW -- 12 months)}

Three parallel tracks:

\begin{enumerate}[label=\textbf{0\alph*:}]
\item \textbf{$D_H/r_d$ sign-change test.} Download DESI DR2 public BAO
  measurements. Plot $D_H/r_d$ residuals vs.\ $z$ relative to best-fit
  $\LCDM$. Look for non-monotonic structure with sign change near
  $z = 1.78 \pm 0.12$. Compare with CPL monotonic prediction.
  \textbf{Personnel}: 1 cosmologist, 2 weeks FTE.
  \textbf{Decision gate}: If sign change detected at $>2\sigma$,
  immediately prioritize the DFD Boltzmann code (hi\_class, \S\ref{sec:boltzmann}).
  If monotonic, reassess DESI tension (may approach falsification).

\item \textbf{GPS 59~ps archival search.} Full W4i10 Proposal~2
  (unchanged). IGS Final clock products, PPP processing, matched
  filter for 59~ps annual signal, multi-GNSS ratio test.
  \textbf{Personnel}: 1 geodesist, 6 months FTE.
  \textbf{Decision gate}: If $\SNR > 5$ and ratio $R = 1.070 \pm 0.01$,
  proceed to multi-GNSS campaign. If null, DFD's $\lambda$-dependent
  prediction is constrained but not killed (the signal scales as
  $(\lambda-1)$).

\item \textbf{LIGO O4 scalar reanalysis.} Updated from W4i10 Proposal~1.
  O4 completed November 2025 with 250 detected mergers. GWTC-4.0
  released August 2025 with 128 new candidates from O4a.
  O4 frequency-dependent squeezing (5.2~dB H1, 6.1~dB L1) improves
  scalar reach to $\lbone \sim 1.5\ee{-14}$.
  \textbf{Key update}: With 250 O4 events (vs.\ 90 in O1--O3), the
  magnetar flare correlation search (Channel~A) has significantly more
  background characterization data. The stochastic scalar background
  search (Channel~C) benefits from $\sim 2.5$ years of O4 data with
  improved sensitivity.
  \textbf{Decision gate}: Upper limit at $1.5\ee{-14}$ improves SQMS
  bound by $10^5\times$. Detection at $>5\sigma$ in 2+ null tests
  would be transformative.
\end{enumerate}

\subsection{Phase I: SQMS (Months 6--30)}

\textbf{SQMS Phase~I} (\$3.2M, 24 months): Unchanged from W4i10
specification. The \$125M DOE renewal (November 2025) provides the
institutional foundation. Primary deliverable: 2$\times$2 cavity
multi-frequency Q-ratio diagnostic.

\begin{itemize}[nosep]
\item Q-ratio prediction: DFD gives 0.49; BCS gives 3.41.
\item Absolute reach at baseline $Q_0 = 5\ee{10}$: $\lbone < 7.8\ee{-10}$.
\item Absolute reach at target $Q_0 = 2\ee{11}$: $\lbone < 1.9\ee{-10}$.
\item Discovery window: \textbf{entire range $1.56\ee{-9}$ to $1.9\ee{-10}$}
  is now open (Dark SRF bound removed).
\end{itemize}

\subsection{Phase II: MEMS + New Proposals (Months 12--36)}

Parallel with Phase~I:

\begin{enumerate}[label=\textbf{II\alph*:}]
\item \textbf{Si$_3$N$_4$ MEMS prototype}: Unchanged from W4i10
  Proposal~3. Q$_{\rm mech} \geq 10^9$ at 20~mK. In-gap
  $\Qpsi$ threshold $2.2\ee{8}$ (4.5$\times$ margin, W4i11).

\item \textbf{Torsion pendulum scalar force search} (NEW, \S\ref{sec:torsion}):
  Exploit the milligram-scale torsion pendulum breakthrough. Place SRF
  cavity psi-source near torsion balance. Independent detection
  geometry from all cavity-based experiments.

\item \textbf{Optical clock diurnal modulation search} (NEW,
  \S\ref{sec:clocks}): Piggyback on CU~Boulder Mt.~Blue Sky campaign.
  Look for diurnal signature in clock comparison residuals.
\end{enumerate}


%=============================================================================
%=============================================================================
\part{New Proposal A: DESI DR2 $D_H/r_d$ Sign-Change Test}
\label{part:desi}
%=============================================================================
%=============================================================================

\begin{tcolorbox}[colback=purple!3!white, colframe=newcolor,
  title=\textbf{NEW PROPOSAL A --- DESI DR2 $D_H/r_d$ SIGN-CHANGE TEST}]
\textbf{Budget}: \$0\\
\textbf{Timeline}: 2 weeks\\
\textbf{Data}: DESI DR2 public BAO measurements (arXiv:2503.14738)\\
\textbf{Personnel}: 1 cosmologist (2 weeks FTE)\\
\textbf{DFD prediction}: Non-monotonic $D_H/r_d$ residual with sign
change at $z = 1.78 \pm 0.12$
\end{tcolorbox}

\section{Scientific Objective}

DFD's dark energy phenomenology does not arise from a modified equation
of state (the CPL parameterization $w_0$-$w_a$ is the \emph{wrong}
observable; W4i11). Instead, DFD produces an ``optical bias'' or
``psi-screen'' that modifies the luminosity distance through the
refractive index $n = e^\psi$ (section02\_formalism.tex, Eq.~2.1).

The key prediction is that $D_H/r_d$ residuals relative to $\LCDM$ are
\emph{non-monotonic}: positive (expansion locally suppressed by
$\mu$-crossover) at $z < z_{\rm cross}$, and negative at $z > z_{\rm cross}$,
with $z_{\rm cross} = 1.78 \pm 0.12$.

CPL models predict \emph{monotonic} departure from $\LCDM$. This
qualitative difference is testable with existing data.

\section{Derivation Chain}

\begin{enumerate}
\item \textbf{Psi-screen mechanism}: The DFD refractive index $n = e^\psi$
  modifies photon geodesics. A photon traversing a psi-gradient
  accumulates an optical path length $\int n\,dl > \int dl$, altering
  the effective Hubble parameter $H(z)$ as inferred from BAO.

\item \textbf{Mu-crossover}: The interpolating function $\mu(x) = x/(1+x)$
  crosses from MOND ($\mu \to x$) to Newtonian ($\mu \to 1$) at
  $x = a/a_0 \sim 1$. In the cosmological context, $x(z) = a_H(z)/a_0$,
  where $a_H = H(z)^2 R_{\rm horizon}$ is the Hubble acceleration.
  The crossover redshift is $z_{\rm cross} \approx 1.78 \pm 0.12$,
  determined entirely by $H_0$, $\Omega_m$, and $a_0$ (zero free
  parameters beyond $\LCDM$ concordance values).

\item \textbf{Sign change}: At $z < z_{\rm cross}$, the universe is in
  the deep-MOND regime for Hubble-scale dynamics, and the psi-screen
  \emph{enhances} $D_H$ relative to $\LCDM$. At $z > z_{\rm cross}$,
  the transition to Newtonian regime \emph{suppresses} $D_H$. The
  residual $\Delta(D_H/r_d) \equiv (D_H/r_d)_{\rm obs} -
  (D_H/r_d)_{\LCDM}$ changes sign.

\item \textbf{CPL contrast}: In CPL parameterization, $w(a) = w_0 + w_a(1-a)$
  is monotonic in $a$ (hence in $z$). The resulting $D_H/r_d$ residual is
  therefore monotonic. Any non-monotonic structure in the residual plot
  \emph{cannot} be captured by CPL and points to non-standard physics.
\end{enumerate}

\section{Method}

\begin{enumerate}[nosep]
\item Download DESI DR2 BAO consensus values for $D_H/r_d$ at each
  effective redshift (7 bins: BGS, LRG1, LRG2, LRG3+ELG1, ELG2,
  QSO, Ly$\alpha$).
\item Adopt the Planck 2018 best-fit $\LCDM$ parameters
  ($H_0 = 67.4$, $\Omega_m = 0.315$) to compute fiducial
  $(D_H/r_d)_{\LCDM}(z_{\rm eff})$.
\item Compute residuals: $\Delta_i = (D_H/r_d)_i^{\rm DESI} -
  (D_H/r_d)_i^{\LCDM}$.
\item Fit a non-monotonic model (DFD psi-screen with $z_{\rm cross}$
  as single parameter) and a monotonic model (CPL with $w_0$, $w_a$).
\item Compute $\Delta\chi^2$ between the two models. Report
  $z_{\rm cross}$ posterior.
\end{enumerate}

\section{Expected Sensitivity}

DESI DR2 achieves 1--2\% precision on $D_H/r_d$ per bin.
The DFD psi-screen predicts $|\Delta(D_H/r_d)/(D_H/r_d)| \sim 1$--3\%
at $z \sim 0.5$ and $z \sim 2.5$. With 7 bins spanning $0.1 < z < 4.2$,
the non-monotonic pattern should be detectable at 2--3$\sigma$ if DFD
is correct.

\textbf{Important caveat}: DESI DR2 analysis (arXiv:2503.14738) reports
a $2.3$--$2.5\sigma$ decreasing trend in $D_H/r_d$ that ``cannot be
explained by physics.'' This trend could be:
(a)~a systematic in DESI BAO or DES SNe,
(b)~real non-standard physics, or
(c)~consistent with the DFD psi-screen (which predicts a
non-monotonic trend that \emph{decreases} at high $z$).
The sign-change test distinguishes (c) from (a) and (b).

\section{Success Criteria}

\begin{itemize}[nosep]
\item \textbf{Detection}: Non-monotonic residual with
  $z_{\rm cross} = 1.78 \pm 0.3$ at $>2\sigma$.
  $\Delta\chi^2(\text{DFD vs CPL}) > 4$.
\item \textbf{Exclusion}: Monotonic residual clearly preferred.
  DFD psi-screen rejected at $>3\sigma$. DESI tension sharpened.
\item \textbf{Inconclusive}: Neither model preferred ($|\Delta\chi^2| < 2$).
  Await DESI DR3 (more redshift bins, smaller error bars).
\end{itemize}


%=============================================================================
%=============================================================================
\part{New Proposal B: Torsion Pendulum Scalar Force Search}
\label{part:torsion}
%=============================================================================
%=============================================================================

\label{sec:torsion}

\begin{tcolorbox}[colback=purple!3!white, colframe=newcolor,
  title=\textbf{NEW PROPOSAL B --- TORSION PENDULUM SCALAR FORCE}]
\textbf{Budget}: \$300--600K\\
\textbf{Timeline}: 18--24 months\\
\textbf{Personnel}: 1 PI (E\"ot-Wash or similar group) + 1 postdoc\\
\textbf{DFD prediction}: Scalar gradient force from SRF cavity psi-source\\
\textbf{Key technology}: Nanofabricated Si$_3$N$_4$ torsion pendulum
with atto-Newton sensitivity
\end{tcolorbox}

\section{Scientific Objective}

DFD predicts that any concentration of EM energy density sources a
scalar perturbation $\dpsi_{\rm EM}$ through Eq.~(\ref{eq:psi_source_main}):
\begin{equation}
\Box\dpsi_{\rm EM} = -\frac{(\lambda-1)\,4\pi G}{c^4}\,u_{\rm EM}.
\label{eq:psi_source_main}
\end{equation}

A test mass at distance $r$ from a spherically symmetric EM source of
total energy $U_{\rm EM}$ experiences a scalar gradient force:
\begin{equation}
F_\psi = \frac{c^2}{2}\,m\,|\nabla\dpsi_{\rm EM}|
= \frac{c^2}{2}\,m\,\frac{(\lambda-1)\,G\,U_{\rm EM}}{c^4\,r^2}
= \frac{(\lambda-1)\,m\,G\,U_{\rm EM}}{2c^2\,r^2}.
\label{eq:F_psi}
\end{equation}

For an SRF cavity with stored energy $U_{\rm EM} = 10$~J (achievable
at $Q_0 = 5\ee{10}$, $E_{\rm acc} = 25$~MV/m in a TESLA 1-cell cavity),
test mass $m = 1$~mg, and distance $r = 1$~cm:
\begin{equation}
F_\psi = \frac{|\lambda-1| \times 10^{-3} \times 6.67\ee{-11} \times 10}
{2 \times (3\ee{8})^2 \times (10^{-2})^2}
= |\lambda-1| \times 3.7\ee{-18}\,\text{N}.
\end{equation}

At the SQMS bound $\lbone = 1.56\ee{-9}$:
\begin{equation}
F_\psi^{\rm max} = 5.8\ee{-27}\,\text{N}.
\end{equation}

\caution{This is 8 orders of magnitude below the atto-Newton
($10^{-18}$~N) sensitivity demonstrated in 2026.}

\subsection{Honest Assessment}

The direct scalar force from an EM-sourced $\dpsi$ is \textbf{too weak}
for current torsion-pendulum technology. This is the same fundamental
limitation identified by the Passive Superiority Theorem (W3i5): the
gravitational coupling $G/c^4$ suppresses all active EM-to-$\psi$
conversion by $\sim 10^{-44}$.

\textbf{However}, the torsion pendulum retains value for two alternative
measurements:

\begin{enumerate}
\item \textbf{Passive Earth-gradient measurement}: The Earth's
  gravitational psi-gradient produces a force on a torsion balance
  that can be compared against Newtonian predictions. DFD predicts
  an anomalous component from the $\mu$-crossover at $a_0$, but in
  the solar neighborhood ($a \gg a_0$), this anomaly is $\sim 10^{-10}$
  of the Newtonian force --- still below current sensitivity.

\item \textbf{MEMS-coupled readout for SQMS Phase~II}: The milligram
  torsion pendulum technology (Si$_3$N$_4$ at 240~$\mu$K) could serve
  as an ultra-sensitive readout for the Phase~II MEMS hybrid detector,
  replacing or supplementing the fiber interferometric readout in the
  W4i10 MEMS proposal.
\end{enumerate}

\textbf{Verdict}: Direct scalar force search KILLED by coupling
strength. Repurpose technology as readout for Phase~II.
Retain as a ``technology development'' proposal at reduced priority.


%=============================================================================
%=============================================================================
\part{New Proposal C: Optical Clock Diurnal Modulation}
\label{part:clocks}
%=============================================================================
%=============================================================================

\label{sec:clocks}

\begin{tcolorbox}[colback=purple!3!white, colframe=newcolor,
  title=\textbf{NEW PROPOSAL C --- OPTICAL CLOCK ALTITUDE COMPARISON}]
\textbf{Budget}: \$0 (piggyback on CU Boulder campaign) or
\$50--200K (dedicated analysis)\\
\textbf{Timeline}: 12 months\\
\textbf{Data}: CU Boulder Mt.~Blue Sky clock comparison
(planned 2025--2026)\\
\textbf{DFD prediction}: Diurnal modulation of clock comparison
residuals tracking solar psi-gradient
\end{tcolorbox}

\section{Scientific Objective}

DFD modifies the gravitational redshift through the psi-screen.
The one-way redshift between two clocks at heights $h_1$ and $h_2$ is:
\begin{equation}
\frac{\Delta f}{f} = \frac{\psi(h_2) - \psi(h_1)}{2}
= \frac{\psi_{\rm grav}(h_2) - \psi_{\rm grav}(h_1)}{2}
+ \frac{\dpsi_{\rm EM}(h_2) - \dpsi_{\rm EM}(h_1)}{2}.
\end{equation}

The first term is the standard GR redshift ($\Delta\Phi/c^2$).
The second term is the DFD anomaly sourced by the Sun's EM luminosity.
The Sun-sourced $\dpsi_{\rm EM}$ varies with the Earth--Sun distance and
the Sun's position relative to the clock pair, producing a \emph{diurnal
modulation} that is absent in GR.

\subsection{Signal estimate}

The Sun's luminosity $L_\odot = 3.83\ee{26}$~W sources a steady-state
$\dpsi_{\rm EM}$ at the Earth's distance:
\begin{equation}
\dpsi_{\rm EM}(r_\oplus) = \frac{(\lambda-1)\,G\,L_\odot}
{c^5\,r_\oplus}
= |\lambda-1| \times 1.9\ee{-26}.
\end{equation}

The \emph{differential} $\dpsi_{\rm EM}$ between 4,350~m and sea level
due to the Sun's gravitational potential gradient:
\begin{equation}
\Delta\dpsi_{\rm EM} \approx \dpsi_{\rm EM} \times
\frac{\Delta h}{r_\oplus} = |\lambda-1| \times 1.3\ee{-32}.
\end{equation}

At $\lbone = 1.56\ee{-9}$: $\Delta f/f \sim 10^{-41}$.

\caution{This is 23 orders of magnitude below the $10^{-18}$
clock precision. The optical clock diurnal modulation test is
\textbf{completely infeasible} at current $\lambda$ bounds.}

\subsection{Revised Role: CLONETS-DS Methodology}

The value of engaging with the CU~Boulder campaign is not the signal
itself but \textbf{establishing the analysis framework} for future
clock network tests:
\begin{itemize}[nosep]
\item Develop the matched-filter template for solar psi-gradient
  diurnal modulation.
\item Characterize systematic noise budget (temperature, atmospheric
  path, fiber phase noise) relevant to CLONETS-DS (European clock
  network, projected 2028).
\item The CLONETS-DS sidereal discriminant (W4i10) requires exactly
  this analysis framework. At continental baselines ($\sim$1000~km),
  the differential psi-gradient from the Earth's gravitational field
  is $\sim 10^{-13}$, which at $\lbone = 10^{-9}$ gives
  $\Delta f/f \sim 10^{-22}$, requiring $\sim 10^6$~s integration.
  Challenging but not impossible with next-generation clocks at
  $10^{-19}$ stability.
\end{itemize}

\textbf{Verdict}: Altitude comparison signal KILLED at current bounds.
Retain as methodology development for CLONETS-DS (2028).


%=============================================================================
%=============================================================================
\part{New Proposal D: Atom Chip Interferometry with Lattice Suspension}
\label{part:atom}
%=============================================================================
%=============================================================================

\begin{tcolorbox}[colback=purple!3!white, colframe=newcolor,
  title=\textbf{NEW PROPOSAL D --- ATOM INTERFEROMETRY SCALAR SEARCH}]
\textbf{Budget}: \$200--500K (analysis + dedicated runs at existing facility)\\
\textbf{Timeline}: 12--18 months\\
\textbf{Facilities}: Stanford AI group, MAGIS-100, or MIGA\\
\textbf{DFD prediction}: Scalar fifth force from Earth's $\psi_{\rm grav}$
gradient, modulated by $\mu$-crossover
\end{tcolorbox}

\section{Scientific Objective}

Lattice atom interferometry (Nature 2024) has achieved unprecedented
precision in measuring gravitational attraction from miniature source
masses. The key advance is optical lattice suspension, which holds atoms
in place during long interrogation times, enabling much higher sensitivity
than free-fall interferometers.

DFD predicts that the gravitational force includes a scalar component
through the $\mu$-function:
\begin{equation}
\mathbf{F}_{\rm DFD} = -m\nabla\Phi_{\rm DFD}
= -\frac{mc^2}{2}\nabla\psi_{\rm grav}
= -\frac{mc^2}{2}\nabla\left[\frac{\psi_{\rm N}}{\mu(|\nabla\psi_{\rm N}|/a_*)}\right],
\end{equation}
where $\psi_{\rm N}$ is the Newtonian potential and $\mu(x) = x/(1+x)$.

For laboratory-scale experiments ($a \gg a_0$, $x \gg 1$, $\mu \approx 1$):
\begin{equation}
\frac{F_{\rm DFD} - F_{\rm Newton}}{F_{\rm Newton}}
\approx \frac{a_0}{a} \sim \frac{10^{-10}}{10}\sim 10^{-11}.
\end{equation}

This $10^{-11}$ fractional force anomaly is \textbf{within reach} of
next-generation atom interferometers. Stanford's apparatus measures
gravitational acceleration with fractional precision $\sim 10^{-9}$
per shot, and lattice interferometry improves this by 2--3 orders of
magnitude.

\subsection{Screening Complication}

\caution{The EFE screening that protects DFD against wide-binary
falsification (W4i11) also suppresses the laboratory signal.}
The solar galactocentric acceleration $a_{\rm ext} = 1.8\,a_0$
screens the MOND anomaly via $\mu_{\rm eff}(x_{\rm eff})$ where
$x_{\rm eff}$ includes the external field. In the solar neighborhood:
\begin{equation}
\frac{\Delta F}{F} \approx \frac{a_0^2}{a \cdot a_{\rm ext}}
\sim \frac{(10^{-10})^2}{10 \times 1.8\ee{-10}} \sim 6\ee{-12}.
\end{equation}

This is reduced by $\sim 2\times$ relative to the unscreened case but
remains within the target sensitivity band.

\subsection{Honest Assessment}

The screened signal ($\sim 6\ee{-12}$ fractional force anomaly) is
at the very edge of projected next-generation atom interferometer
sensitivity. Moreover:
\begin{itemize}[nosep]
\item Current lattice interferometers achieve $\sim 10^{-9}$ per shot.
\item With $10^4$ shots and systematic control: $\sim 10^{-11}$.
\item The DFD signal ($6\ee{-12}$) requires systematic control at
  the $10^{-12}$ level, which is not yet demonstrated.
\item Source mass characterization (knowing $M_{\rm source}$ to better
  than $10^{-12}$ fractionally) is extremely challenging.
\end{itemize}

\textbf{Verdict}: Promising for 2028--2030 timeframe with next-generation
facilities (MAGIS-100, MIGA). Not actionable in 2026 with current
sensitivity. \textbf{Retain as medium-term proposal.} The key deliverable
in 2026 is a feasibility study estimating systematic budgets.


%=============================================================================
%=============================================================================
\part{New Proposal E: Muon Storage Ring Anomalous Precession}
\label{part:muon}
%=============================================================================
%=============================================================================

\begin{tcolorbox}[colback=purple!3!white, colframe=newcolor,
  title=\textbf{NEW PROPOSAL E --- MUON $g{-}2$ STORAGE RING REINTERPRETATION}]
\textbf{Budget}: \$0 (reanalysis of existing data)\\
\textbf{Timeline}: 3--6 months\\
\textbf{Data}: Fermilab $g{-}2$ Run 1--6 data\\
\textbf{DFD prediction}: Null ($\Delta a_\mu^{\rm DFD} \sim 10^{-15}$,
4 orders below sensitivity)
\end{tcolorbox}

\section{Scientific Objective}

The Fermilab muon $g{-}2$ experiment measures the anomalous magnetic
moment of the muon with extraordinary precision. DFD predicts a null
result: the psi-field correction to $g{-}2$ is
$\Delta a_\mu^{\rm DFD} \sim |\lambda-1| \times G m_\mu^2/(\hbar c)
\sim 10^{-15}$ (MASTER\_FINDINGS\_CORPUS), which is 4 orders of magnitude
below the current experimental precision of $\sim 10^{-11}$.

\textbf{However}, the $g{-}2$ storage ring magnetic field ($B = 1.45$~T,
stored in a $\sim 14$~m radius ring) represents a significant EM energy
density source. The ring field energy is
$U_B = B^2 V/(2\mu_0) \sim 10^5$~J. This sources a $\dpsi_{\rm EM}$:
\begin{equation}
\dpsi_{\rm EM}^{\rm ring} \sim \frac{|\lambda-1| \cdot G \cdot U_B}
{c^4 \cdot R_{\rm ring}} \sim |\lambda-1| \times 8\ee{-36}.
\end{equation}

This is negligible. \textbf{Verdict}: No DFD signal in muon $g{-}2$.
Published as a null-prediction paper (establishing that DFD does
\emph{not} contribute to the $g{-}2$ anomaly, in contrast to some
BSM scenarios).


%=============================================================================
%=============================================================================
\part{New Proposal F: Precision Casimir Measurement as Scalar Probe}
\label{part:casimir}
%=============================================================================
%=============================================================================

\begin{tcolorbox}[colback=purple!3!white, colframe=newcolor,
  title=\textbf{NEW PROPOSAL F --- PRECISION CASIMIR SCALAR PROBE}]
\textbf{Budget}: \$100--300K (dedicated analysis + systematic characterization)\\
\textbf{Timeline}: 12 months\\
\textbf{Facilities}: Existing Casimir measurement groups (Yale, IUPUI, Purdue)\\
\textbf{DFD prediction}: Modified Casimir force via
$n = e^\psi$ refractive index
\end{tcolorbox}

\section{Scientific Objective}

The Casimir effect arises from quantum vacuum fluctuations of the
EM field between conducting surfaces. In DFD, the vacuum has an
effective refractive index $n = e^\psi$ that modifies the mode
spectrum and hence the Casimir force.

The DFD-modified Casimir pressure between parallel plates of area $A$
separated by distance $d$:
\begin{equation}
P_{\rm Casimir}^{\rm DFD} = -\frac{\pi^2 \hbar c}{240\,d^4}
\times e^{-\psi_{\rm bg}} \approx P_{\rm Casimir}^{\rm QED}
\times (1 - \psi_{\rm bg}),
\end{equation}
where $\psi_{\rm bg}$ is the gravitational potential at the lab location
($\psi_{\rm bg} \sim GM_\oplus/(c^2 R_\oplus) \sim 7\ee{-10}$).

The fractional DFD anomaly in the Casimir force is:
\begin{equation}
\frac{\Delta F_{\rm Casimir}}{F_{\rm Casimir}} \sim \psi_{\rm bg}
\sim 7\ee{-10}.
\end{equation}

\textbf{This is comparable to the current precision frontier of Casimir
measurements} ($\sim 0.1$--1\% or $10^{-3}$--$10^{-2}$). However:

\begin{itemize}[nosep]
\item The $\psi_{\rm bg}$ correction is a constant offset, not a
  modulated signal. It cannot be distinguished from a systematic
  calibration error in the plate separation or dielectric function.
\item The prediction is for $\psi_{\rm grav}$, not $\dpsi_{\rm EM}$.
  Since $\psi_{\rm grav} = -\Phi/c^2$ (where $\Phi$ is the Newtonian
  potential), this is just the standard weak-field gravitational
  redshift of the Casimir force --- already included in GR.
\item The \emph{DFD-specific} anomaly (from $\lambda \neq 1$) would
  require an EM-sourced $\dpsi_{\rm EM}$ gradient between the plates,
  which is negligible at $\sim 10^{-36}$.
\end{itemize}

\textbf{Verdict}: KILLED. The measurable Casimir correction ($7\ee{-10}$)
is gravitational (GR-standard), not DFD-specific. The DFD-specific
EM-sourced correction is $\sim 10^{-36}$, completely unmeasurable.


%=============================================================================
%=============================================================================
\part{The DFD Boltzmann Code: Critical Infrastructure}
\label{part:boltzmann}
%=============================================================================
%=============================================================================

\label{sec:boltzmann}

\begin{tcolorbox}[colback=red!5!white, colframe=warncolor,
  title=\textbf{URGENT: DFD BOLTZMANN CODE NEEDED FOR ALL COSMOLOGICAL TESTS}]
The DFD Boltzmann code (hi\_class modification) is the single most
important piece of \textbf{infrastructure} for the experimental programme.
Without it:
\begin{itemize}[nosep]
\item The $D_H/r_d$ sign-change test (Proposal~A) can only be done
  approximately using analytical psi-screen templates.
\item The DESI tension (3.1--3.5$\sigma$, W4i11) cannot be properly
  assessed.
\item The Euclid $S_8$ scale-dependence test cannot generate
  quantitative predictions.
\item The neutrino mass sum prediction ($61.4$~meV) cannot be confronted
  with CMB-S4 + DESI joint constraints.
\end{itemize}
\textbf{Estimated effort}: 2,520 lines of code, 4.5 months for a
competent cosmologist (W4i12 Boltzmann Spec). Maps onto Horndeski
$\alpha$-functions ($\alpha_M = d\psi_0/(aH)$, $\alpha_B = \alpha_T = 0$).
\textbf{Budget}: \$80--150K (1 postdoc, 6 months).
\end{tcolorbox}


%=============================================================================
%=============================================================================
\part{Consolidated Experimental Programme}
\label{part:consolidated}
%=============================================================================
%=============================================================================

\section{Complete Budget Summary}

\begin{center}
\renewcommand{\arraystretch}{1.3}
\begin{longtable}{@{}clcccc@{}}
\toprule
\textbf{Rank} & \textbf{Experiment} & \textbf{Cost} & \textbf{Time} &
  \textbf{Status} & \textbf{Source} \\
\midrule
\endhead
0a & $D_H/r_d$ sign change & \$0 & 2 wks & NEW (W4i13) & Proposal A \\
0b & GPS 59~ps archival & \$50--100K & 6 mo & W4i10 P2 & Unchanged \\
0c & LIGO O4 scalar & \$100--500K & 12 mo & W4i10 P1 & Updated \\
0d & Rotation curve reanalysis & \$0 & 3--6 mo & W4i10 P4 & Unchanged \\
\midrule
$\infty$ & DFD Boltzmann code & \$80--150K & 6 mo & URGENT & W4i12 Spec \\
\midrule
I & SQMS Phase~I & \$3.2M & 24 mo & W4i10 P11 & Wider window \\
II & Si$_3$N$_4$ MEMS & \$200--500K & 18 mo & W4i10 P3 & Unchanged \\
\midrule
N1 & Torsion pendulum readout & \$300--600K & 24 mo & NEW (W4i13) &
  Technology dev \\
N2 & Atom interferometry study & \$50--100K & 6 mo & NEW (W4i13) &
  Feasibility \\
\midrule
\multicolumn{2}{l}{\textbf{KILLED proposals:}} & & & & \\
--- & Torsion direct force & --- & --- & KILLED & $G/c^4$ too weak \\
--- & Optical clock altitude & --- & --- & KILLED & Signal $10^{-41}$ \\
--- & Neutron interferometry & --- & --- & KILLED & Signal $10^{-22}$ rad \\
--- & Casimir scalar probe & --- & --- & KILLED & GR-standard, not DFD \\
--- & Muon $g{-}2$ reinterpretation & --- & --- & NULL & $10^{-15}$, publish \\
--- & Dark SRF reinterpretation & --- & --- & INVALIDATED & W4i12 \\
\bottomrule
\end{longtable}
\end{center}

\textbf{Total Phase 0 investment}: \$150--750K (archival + Boltzmann code).

\textbf{Total Phase I investment}: \$3.2M (SQMS, unchanged).

\textbf{Total new W4i13 proposals}: \$350--700K (torsion readout + atom
interferometry feasibility).


\section{Five Killed Proposals and Why}

The Passive Superiority Theorem (W3i5) dominates the landscape of
``new'' experimental ideas:

\begin{center}
\renewcommand{\arraystretch}{1.3}
\begin{tabular}{@{}p{4cm}p{3cm}p{6cm}@{}}
\toprule
\textbf{Proposal} & \textbf{Predicted Signal} & \textbf{Kill Mechanism} \\
\midrule
Torsion direct force & $5.8\ee{-27}$~N & $G/c^4$ coupling: 8 orders
  below atto-Newton sensitivity \\
Optical clock altitude & $\Delta f/f \sim 10^{-41}$ & $G/c^4$ coupling:
  23 orders below clock precision \\
Neutron interferometry & $10^{-22}$~rad & $G/c^4$ coupling:
  18 orders below phase sensitivity \\
Casimir scalar probe & $\sim 10^{-36}$ & DFD-specific correction
  negligible; gravitational part is standard GR \\
Dark SRF LSW & INVALIDATED & Spin-0 vs spin-1 conflation;
  $G/c^4$ active conversion too weak \\
\bottomrule
\end{tabular}
\end{center}

\textbf{Common theme}: Every experiment that relies on \emph{active}
EM-to-$\psi$ conversion (creating a $\dpsi_{\rm EM}$ signal from a
laboratory EM source) is killed by the $G/c^4 = 9.3\ee{-45}$ coupling
constant. This is the Passive Superiority Theorem in action: passive
observation of existing gravitational $\psi_{\rm grav}$ gradients
(Earth, Sun, cosmos) beats active EM sourcing by 14+ orders of magnitude.

\textbf{The surviving experiments} are all passive:
\begin{itemize}[nosep]
\item DESI: Observes cosmological psi-screen (passive).
\item GPS: Observes Sun-sourced psi-gradient (passive, astrophysical source).
\item LIGO: Observes astrophysical psi-perturbations (passive).
\item SQMS: Measures psi back-reaction on cavity Q (passive diagnostic).
\item MEMS: Measures in-gap $\Qpsi$ (passive threshold test).
\item Rotation curves: Observes galactic dynamics (passive).
\end{itemize}


\section{Decision Tree}

\begin{center}
\fbox{\parbox{0.9\textwidth}{
\textbf{Step 1} (NOW): Run $D_H/r_d$ sign-change test.\\
$\to$ If sign change detected ($>2\sigma$): Fast-track Boltzmann code.
DFD cosmological sector strengthened.\\
$\to$ If monotonic ($>3\sigma$): DESI tension $>4\sigma$.
DFD cosmological sector under severe pressure.
Technology patents unaffected.\\[6pt]
\textbf{Step 2} (Months 1--6): GPS 59~ps archival search.\\
$\to$ If detected: Smoking gun for DFD ($\lambda \neq 1$).
Immediate multi-GNSS campaign.\\
$\to$ If null: Constrains $\lambda$ but does not kill DFD.\\[6pt]
\textbf{Step 3} (Months 6--30): SQMS Phase~I.\\
$\to$ Q-ratio = 0.49: DFD confirmed. All technology patents activated.\\
$\to$ Q-ratio = 3.41: BCS standard. $\lbone$ pushed to $<1.9\ee{-10}$.\\
$\to$ Intermediate: Interesting physics, needs interpretation.\\[6pt]
\textbf{Step 4} (Months 12--36): LIGO O4 + MEMS in parallel.\\
$\to$ LIGO detection: $\lbone < 1.5\ee{-14}$, transformative.\\
$\to$ MEMS threshold crossed: Phase~II detector justified.\\[6pt]
\textbf{Step 5} (2028+): CLONETS-DS, Euclid, JUNO, CMB-S4.\\
$\to$ Decisive multi-front confrontation.
}}
\end{center}


%=============================================================================
\section{Inconsistencies and Open Issues}
\label{sec:inconsistencies}
%=============================================================================

\begin{enumerate}
\item \textbf{DESI DR2 $D_H/r_d$ trend}: The 2.3--2.5$\sigma$ decreasing
  trend reported in DESI DR2 cross-checks (arXiv:2510.04179) is
  potentially consistent with DFD's psi-screen, but the DESI collaboration
  attributes it to systematics. A DFD-specific analysis is needed to
  distinguish the two interpretations. \textbf{Action}: This is exactly
  what the Boltzmann code enables.

\item \textbf{Passive Superiority vs.\ technology patents}: The Passive
  Superiority Theorem kills all active EM-to-$\psi$ conversion experiments.
  But the technology patents (P2--P4, P6--P7, P9--P13, P15) \emph{require}
  efficient EM-to-$\psi$ conversion. The resolution is that technology
  applications use resonant cavity buildup at cooperativity $C \geq 1$
  (Fundamental Conversion Limit Theorem: $\eta = 4C/(1+C)^2 \to 1$
  at $C = 1$). The detection programme uses passive observation precisely
  because it avoids the $G/c^4$ suppression. These are not contradictory:
  the technology operates at $C \geq 1$ (designed), while laboratory
  detection operates at $C \ll 1$ (limited by current $Q$ and $\lambda$
  bounds). \textbf{No inconsistency.}

\item \textbf{SQMS Phase~I timing vs.\ DESI urgency}: If the $D_H/r_d$
  test shows monotonic residuals ($>3\sigma$), the DESI tension rises
  above $4\sigma$ and the cosmological sector of DFD faces potential
  falsification. Should SQMS Phase~I proceed regardless?
  \textbf{Answer}: Yes. The technology patents depend on $\lambda$,
  not on the dark energy model. Even if DFD's cosmological sector
  is falsified, the detection of $\lambda \neq 1$ via Q-ratio
  diagnostic would validate the EM-$\psi$ coupling and all
  technology applications.

\item \textbf{Atom interferometry screening}: The EFE screening that
  saves DFD from wide-binary falsification (W4i11) simultaneously
  suppresses the laboratory atom interferometry signal by $\sim 2\times$.
  The screened signal ($6\ee{-12}$ fractional anomaly) requires
  systematic control at a level not yet demonstrated. This is
  a genuine tension between ``DFD is not falsified by wide binaries''
  and ``DFD is testable in the lab.'' \textbf{The resolution is
  temporal}: by 2028--2030, MAGIS-100 and MIGA should achieve the
  necessary sensitivity.
\end{enumerate}


%=============================================================================
\section*{Summary of Derivation Chains}
%=============================================================================

\begin{enumerate}
\item \textbf{$D_H/r_d$ sign change}: Action (Eq.~2.6, section02)
  $\to$ $\mu(x) = x/(1+x)$ $\to$ cosmological $\mu$-crossover at
  $z_{\rm cross} = 1.78$ $\to$ non-monotonic $D_H/r_d$ residuals $\to$
  testable with DESI DR2 public data.

\item \textbf{Torsion pendulum force}: Source equation
  (Eq.~\ref{eq:psi_source_main}) $\to$ $\dpsi_{\rm EM} \propto
  (\lambda-1) G U_{\rm EM}/(c^4 r)$ $\to$ force
  $F_\psi = (c^2/2) m |\nabla\dpsi|$ $\to$ $5.8\ee{-27}$~N at
  SQMS bound $\to$ KILLED (8 orders below sensitivity).

\item \textbf{SQMS Phase~I (recalibrated)}: Passive Q-measurement $\to$
  $\psi$-channel loss rate $\propto |\lambda-1|^2 / Q_\psi$ $\to$
  Q-ratio diagnostic 0.49 vs 3.41 $\to$ discovery window
  $[1.56\ee{-9}, 1.9\ee{-10}]$ now fully open.

\item \textbf{Atom interferometry}: $\mu$-crossover in lab gravity
  $\to$ fractional anomaly $a_0/a \sim 10^{-11}$ $\to$ EFE screening
  reduces to $6\ee{-12}$ $\to$ within projected 2028--2030 sensitivity.

\item \textbf{Killed proposals}: All share the derivation chain
  EM source $\to$ $\dpsi_{\rm EM} \propto G/c^4$ $\to$ signal
  $\sim 10^{-27}$ to $10^{-41}$ $\to$ far below any current or
  planned detector.
\end{enumerate}

\end{document}
